% !TEX TS-program = XeLaTeX
% use the following command:
% all document files must be coded in UTF-8
\documentclass[english]{textolivre}
% build HTML with: make4ht -e build.lua -c textolivre.cfg -x -u article "fn-in,svg,pic-align"

\journalname{Texto Livre}
\thevolume{19}
%\thenumber{1} % old template
\theyear{2026}
\receiveddate{\DTMdisplaydate{2025}{9}{15}{-1}} % YYYY MM DD
\accepteddate{\DTMdisplaydate{2026}{2}{12}{-1}}
\publisheddate{\DTMdisplaydate{2026}{5}{27}{-1}}
\corrauthor{Carlos Henrique R. Valadares}
\articledoi{10.1590/1983-3652.2026.61768}
%\articleid{NNNN} % if the article ID is not the last 5 numbers of its DOI, provide it using \articleid{} commmand 
% list of available sesscions in the journal: articles, dossier, reports, essays, reviews, interviews, editorial
\articlesessionname{articles}
\runningauthor{Valadares} 
%\editorname{Leonardo Araújo} % old template
\sectioneditorname{Daniervelin Pereira~\orcid{0000-0003-1861-3609}}
\layouteditorname{Leonardo Araújo~\orcid{0000-0003-3884-2177}}

\title{Ludic possibilities: exploring the affordances of games in additional languages education}
\othertitle{Possibilidades lúdicas: explorando as \textit{affordances} de jogos na educação de línguas adicionais}
% if there is a third language title, add here:
%\othertitle{Artikelvorlage zur Einreichung beim Texto Livre Journal}

\author[1]{Carlos Henrique Rodrigues Valadares~\orcid{0000-0002-8999-6216}\thanks{Email: \href{mailto:prof.carlosvaladares@gmail.com}{prof.carlosvaladares@gmail.com}}}
\affil[1]{Universidade Federal de Minas Gerais, Faculdade de Letras, Belo Horizonte, MG, Brasil.}

\addbibresource{article.bib}
% use biber instead of bibtex
% $ biber article

% used to create dummy text for the template file
\definecolor{dark-gray}{gray}{0.35} % color used to display dummy texts
\usepackage{lipsum}
\SetLipsumParListSurrounders{\colorlet{oldcolor}{.}\color{dark-gray}}{\color{oldcolor}}

% used here only to provide the XeLaTeX and BibTeX logos
\usepackage{hologo}

% if you use multirows in a table, include the multirow package
\usepackage{multirow}

% provides sidewaysfigure environment
\usepackage{rotating}

% CUSTOM EPIGRAPH - BEGIN 
%%% https://tex.stackexchange.com/questions/193178/specific-epigraph-style
\usepackage{epigraph}
\renewcommand\textflush{flushright}
\makeatletter
\newlength\epitextskip
\pretocmd{\@epitext}{\em}{}{}
\apptocmd{\@epitext}{\em}{}{}
\patchcmd{\epigraph}{\@epitext{#1}\\}{\@epitext{#1}\\[\epitextskip]}{}{}
\makeatother
\setlength\epigraphrule{0pt}
\setlength\epitextskip{0.5ex}
\setlength\epigraphwidth{.7\textwidth}
% CUSTOM EPIGRAPH - END

% to use IPA symbols in unicode add
%\usepackage{fontspec}
%\newfontfamily\ipafont{CMU Serif}
%\newcommand{\ipa}[1]{{\ipafont #1}}
% and in the text you may use the \ipa{...} command passing the symbols in unicode

% LANGUAGE - BEGIN
% ARABIC
% for languages that use special fonts, you must provide the typeface that will be used
% \setotherlanguage{arabic}
% \newfontfamily\arabicfont[Script=Arabic]{Amiri}
% \newfontfamily\arabicfontsf[Script=Arabic]{Amiri}
% \newfontfamily\arabicfonttt[Script=Arabic]{Amiri}
%
% in the article, to add arabic text use: \textlang{arabic}{ ... }
%
% RUSSIAN
% for russian text we also need to define fonts with support for Cyrillic script
% \usepackage{fontspec}
% \setotherlanguage{russian}
% \newfontfamily\cyrillicfont{Times New Roman}
% \newfontfamily\cyrillicfontsf{Times New Roman}[Script=Cyrillic]
% \newfontfamily\cyrillicfonttt{Times New Roman}[Script=Cyrillic]
%
% in the text use \begin{russian} ... \end{russian}
% LANGUAGE - END

% EMOJIS - BEGIN
% to use emoticons in your manuscript
% https://stackoverflow.com/questions/190145/how-to-insert-emoticons-in-latex/57076064
% using font Symbola, which has full support
% the font may be downloaded at:
% https://dn-works.com/ufas/
% add to preamble:
% \newfontfamily\Symbola{Symbola}
% in the text use:
% {\Symbola }
% EMOJIS - END

% LABEL REFERENCE TO DESCRIPTIVE LIST - BEGIN
% reference itens in a descriptive list using their labels instead of numbers
% insert the code below in the preambule:
%\makeatletter
%\let\orgdescriptionlabel\descriptionlabel
%\renewcommand*{\descriptionlabel}[1]{%
%  \let\orglabel\label
%  \let\label\@gobble
%  \phantomsection
%  \edef\@currentlabel{#1\unskip}%
%  \let\label\orglabel
%  \orgdescriptionlabel{#1}%
%}
%\makeatother
%
% in your document, use as illustraded here:
%\begin{description}
%  \item[first\label{itm1}] this is only an example;
%  % ...  add more items
%\end{description}
% LABEL REFERENCE TO DESCRIPTIVE LIST - END


% add line numbers for submission
%\usepackage{lineno}
%\linenumbers

\begin{document}
\maketitle

\begin{polyabstract}
\begin{abstract}
This study\footnote{This study is derived from the author’s Master’s dissertation titled Tomando o controle: affordances de jogos digitais para aprender e ensinar línguas adicionais, submitted to Programa de Pós-Graduação em Estudos Linguísticos at Universidade Federal de Minas Gerais (POSLIN/UFMG) in 2023. Advisor: Prof. Ronaldo Corrêa Gomes Junior.} investigates the use of digital games in additional language learning, focusing on teachers’ perceptions of digital game affordances for additional language teaching. By conducting a five-week workshop on teaching with digital games, it sought to understand how teachers integrate digital games into their practices and analyze the teaching materials they created. Data comprised lesson plans, feedback forms, and narrative reports. The study identified four affordance categories: interact, listen and read, create, and perform tasks, along with ten subcategories. The research emphasizes the importance of teachers as mediators in using games for additional language teaching, and the recognition of teachers’ central role in leveraging digital technologies for effective learning.

\keywords{Teacher education \sep Affordances \sep Digital game-based language learning \sep Ludic language pedagogy}
\end{abstract}

\begin{portuguese}
\begin{abstract}
Este estudo investiga o uso de jogos digitais na aprendizagem de línguas adicionais, com foco nas percepções de professores sobre as \textit{affordances} dos jogos digitais para o ensino de línguas adicionais. Por meio da realização de uma oficina de cinco semanas sobre ensino com jogos digitais, buscou-se compreender como os professores integram jogos digitais em suas práticas e analisar os materiais de ensino que eles criaram. Os dados incluíram planos de aula, formulários de \textit{feedback} e relatórios narrativos. O estudo identificou quatro categorias de possibilidades: interagir, ouvir e ler, criar e realizar tarefas, juntamente com dez subcategorias. A pesquisa enfatiza a importância dos professores como mediadores no uso de jogos para o ensino de línguas adicionais, e o reconhecimento do papel central dos professores na alavancagem de tecnologias digitais para uma aprendizagem eficaz.

\keywords{Formação de professores \sep Affordances \sep Aprendizagem baseada em jogos \sep Pedagogia lúdica da linguagem}
\end{abstract}
\end{portuguese}
% if there is another abstract, insert it here using the same scheme
\end{polyabstract}

\section{Introduction}\label{sec-intro}
Digital games have become an integral part of daily life for many individuals, spanning various age groups and serving purposes ranging from leisure to professional endeavors. This widespread adoption has also permeated the educational landscape, where the use of digital information and communication technologies (DICTs) has gained significant traction. According to \textcite[p. 45]{kenski2007educacao}, “the choice of a particular type of technology profoundly alters the nature of the educational process and communication between participants.”

This study revolves around the use of digital games as a tool for language teaching and learning. Specifically, we aim to explore how teachers of additional languages (both in-service and in-training contexts) perceive and utilize digital games in their pedagogical practices. Despite the growing body of research on digital games in education, there is a notable gap in studies that move beyond theoretical discussions to practical classroom applications and teacher training programs \cite{york2021playing}. Within this framework, digital technologies are regarded as tools capable of expanding learning opportunities. The ecological perspective \cite{gibson1986ecological,vanlier2000input,vanlier2004ecology,vanlier2010ecology} underscores the classroom as a plural space where diverse realities converge, enabling the sharing of knowledge that prioritizes the context of students’ lives.

Existing research has demonstrated the potential of digital games in language learning. For instance, \textcite{gee2003videogames} argues that the immersive experiences provided by games can facilitate learning, not only in language but across various fields of knowledge. Similarly, studies by \textcite{rankin2006rpg,rama2012wow} have shown that games like Everquest and World of Warcraft offer rich environments for language practice, particularly through social interactions and collaborative tasks. However, there is a need for more research that focuses on how teachers can harness these affordances in their teaching practices, especially in additional language learning contexts.

The main goal of this research was to analyze digital game affordances perceived by additional language teachers. By doing so, we seek to provide a framework that not only enhances language learning but also empowers teachers to design and implement game-based activities that align with their curriculum and their students’ needs. This approach is grounded in the theoretical-methodological framework of “Ludic Language Pedagogy” (LLP) \cite{york2021playing}, which emphasizes the teacher’s role as a mediator between games and language teaching.

\section{Research context}\label{sec-normas}
This study was conducted within the context of a continuing education workshop for both in-service and in-training (Modern Languages undergraduates) additional language teachers. The workshop, titled “Play, Design \& Teach,” was conducted online over five weeks, using the virtual platform Google Classroom to accommodate participants from diverse geographical locations in Brazil. The six participants, aged between 22 and 36 years, were primarily English language teachers, with one participant teaching Spanish. All participants had former experience in teaching additional languages, and most were pursuing higher education degrees, either at the undergraduate, master, or doctoral level. To ensure the integrity of the research and the protection of those involved, the study followed established ethical protocols and received official institutional approval\footnote{This study is linked to the research project “Fora da Escola: Histórias de Aprendizagem Informal Online de Inglês,” coordinated by Professor Ronaldo Corrêa Gomes Junior and approved by the Research Ethics Committee of the Federal University of Minas Gerais (CAAE: 88447718.7.0000.5149).} (\Cref{tab1}).

\begin{table}[h!]
\begin{threeparttable}
\caption{Workshop participants’ profile}
\label{tab1}
\small
\centering
\begin{tabular}{p{1cm} p{1cm} p{1cm} p{2,5cm} p{3cm} p{3,5cm}}
\hline
Name & Gender & Age & Educational Level & Teaches & Teaching settings \\
\hline
Antúrio & Nb & 33 & Undergraduate (ongoing) &
English, German, Portuguese for foreigners & Language school, Private classes \\

Girassol & M & 23 & Undergraduate (ongoing) &
English, Portuguese for foreigners & Extension course \\

Gérbera & F & 36 & Master’s &
English & Public school, Private classes, Extension course \\

Lavanda & F & 22 & Undergraduate (ongoing) &
English, Brazilian Sign Language (Libras) & Language school, Public university/college \\

Manacá & M & 25 & Undergraduate (ongoing) &
English & Language school \\

Sálvia & F & 30 & PhD &
Spanish & Private school, Private classes, Public university/college \\
\hline
\end{tabular}

\vspace{0.3cm}

\footnotesize
\begin{tablenotes}
\small{
\item{Nb = non-binary; M = male; F = female. Pseudonyms are common Brazilian plant names.}
 }
 \source{The author.}
\end{tablenotes}
\end{threeparttable}
\end{table}

The learners targeted by the participants in their teaching contexts varied in age and proficiency levels, ranging from young learners to adults, and from beginner to advanced additional language proficiency. The participants taught in diverse educational settings, including public and private schools, language institutes, and higher education institutions. The five-week course aimed to equip these teachers with the skills and knowledge to integrate digital games into their language teaching practices, regardless of the specific age or proficiency level of their students.

Access to technology among the participants and their students was a key consideration in the workshop planning. All participants reported having access to computers and smartphones, but with a diverse range of technological knowledge levels. Despite these disparities, the course was designed to be inclusive, offering activities that could be adapted to different levels of technological access. For instance, some activities required only a smartphone and an internet connection, while others could be conducted using free or low-cost digital games available on various platforms.

The teaching plan and syllabi for the course were designed to encourage reading and reflection on texts and videos about gaming, as well as the planning and development of activities involving the use of digital games for language education. As digital games have different genres and platforms, distinct methodological approaches were used in each activity.

The project title, “Play, Design \& Teach,” aimed to encapsulate the three primary components of the course: \textbf{Play} was a call to action for new, ludic experiences. The term ‘Play’ embodies a dual meaning, encompassing both engaging with games and interacting with media such as videos and music. Throughout the workshop, activities were explored based on the different levels of interaction proposed in LLP \cite{york2021playing}. Thus, in addition to digital games (activities \textit{within} the game), this section contained gameplay videos (activities \textit{around} the game) and research materials with games (\textit{ecological framing} of the game), seeking a theoretical and practical vision that would enable participants to understand the use of digital games as potential teaching tools.

\textbf{Design} refers to production and planning. Whether lessons or activities, designing is part of an educator’s job. In addition, design refers to one of the facets of the Pedagogy of Multiliteracies, where active participation is central to the learning process \cite{newlondongroup1996multiliteracies,cope2015things}. In the \textbf{Design} section, participants created lesson plans, activities, exercises, and infographics, planning on how to integrate the content seen in \textbf{Play} to the reality of their classrooms.

The \textbf{Teach} section was a way to discuss teaching practices. Here, participants were encouraged to share their experiences with the activities and materials produced, discuss their potential, and ask questions. \textbf{Teach} was, therefore, a space for sharing. Participant reports from this section enabled an analysis of perceived affordances, mediated by targeted questions designed to elicit them.

\section{Workshop description and procedures}\label{sec-conduta}
As part of the planned dynamics for the \textbf{Play} section, participants were led to explore games related to a weekly topic. As illustrated in \Cref{tab2}, there were four weeks of interaction with games, while the last week, “autonomy,” focused on reflecting on ways to motivate students to proactively learn an additional language through digital games. The “planned games” column contains games mentioned in the activity plans for each week, whereas “games selected by the participants” lists the games that each participant utilized for carrying out the activities, which may or may not be the digital games presented during the workshop planning.

\begin{table}[h!]
\begin{threeparttable}
\caption{Digital games used in the “Play, Design \& Teach” workshop}
\label{tab2}
\centering
\small
\begin{tabular}{p{1cm}p{2.5cm}p{5cm}p{5cm}}
\hline
Week & Theme & Planned Games & Games selected by the participants \\
\hline

1 & Simulation &
The Sims &
The Sims 4, The Sims Mobile \\

2 & Party Games &
Among Us, Gartic Phone &
Spyfall, Among Us, Gartic Phone \\

3 & Singleplayer &
Arida: Backlands Awakening, Child of Light, A Plague Tale: Innocence &
Arida: Backlands Awakening, The Stanley Parable, Stray, Stardew Valley \\

4 & Sandbox &
Minecraft &
Minecraft \\

5 & Autonomy &
Instructional video &
-- \\
\hline
\end{tabular}

\vspace{0.3cm}

\footnotesize
\source{The author.}
\end{threeparttable}
\end{table}

The first and fourth weeks focused on specific digital games, The Sims and Minecraft, respectively. During these weeks, participants played the games and then designed an activity (in The Sims) and remixed existing lesson plans (in Minecraft) for their own class contexts. Given their distinct genres (Simulation and Sandbox), practical accessibility via computers and smartphones, and extensive online communities, The Sims and Minecraft were selected for more specific activities.

The second week shifted its focus to Party Games, which are typically multiplayer and emphasize social interaction \cite{oliveira2013videogames}. Games like Among Us, Gartic Phone, and Spyfall were introduced as examples.

The third week centered on single-player games. Participants were encouraged to watch a gameplay video, since the time-consuming nature of running more complex digital games (hardware or console requirements, install and execution time) could prevent them from carrying on this activity. Gameplay videos align with different levels of interaction with and about games in LLP \cite{york2021playing}, and were selected considering various languages, audiences, and cultures. Despite the provided videos, learners were encouraged to search within their own libraries for the games they most identified with to produce the activities. This allowed for creative combinations and proposals from each individual participant. Both the second and third weeks emphasized independent research, exploration, and activity development by the participants.

The final week focused on autonomy for language learning, encouraging participants to reflect on how digital games could promote self-directed learning among their students. Participants were asked to develop strategies for motivating students to use digital games proactively for language learning outside the classroom. This week did not focus on a specific game, but rather on the broader concept of leveraging games to foster learner autonomy.

Throughout the workshop, participants were actively involved in designing and sharing their lesson plans and activities. They provided feedback on each other’s ideas and discussed the potential challenges and benefits of using digital games in their teaching contexts. The participants’ input was crucial in shaping the final outcomes of the workshop, as they brought diverse perspectives and experiences to the table. For example, Sálvia proposed an activity using The Sims to teach physical and psychological characteristics in Spanish, where students would create characters and describe them in the target language. This activity highlighted the potential of digital games to integrate multiple language skills, such as vocabulary, grammar, and descriptive writing.

\section{Results and analysis}\label{sec-fmt-manuscrito}
The table of game affordances was developed through a systematic analysis of the data collected during the “Play, Design \& Teach” workshop. The data included participants’ lesson plans, weekly feedback forms, forum interactions, and narrative reports.

The analysis followed a qualitative approach, in which the affordances were characterized based on the opportunities for action present in the digital games used. As a result, only the affordances that, in some way, created direct relationships with additional language teaching and learning were selected to compose the categories of this research. These categories were named using verbs that indicate player actions. The coding process was facilitated by the ATLAS.ti 22 software, which allowed for the organization and classification of the data into meaningful categories.

The affordances were grouped into four main categories, each representing a different type of interaction or opportunity for language learning within the context of digital games. These categories were further divided into subcategories to capture the nuances of how these affordances manifest in different games, as shown in \Cref{tab3}.

\begin{table}[h!]
\begin{threeparttable}
\caption{Categorization of game affordances extracted from the teachers’ reports}
\label{tab3}
\centering
\small
\begin{tabular}{p{3cm}p{4,5cm}p{5,5cm}}
\hline
Affordance & Subcategory & Related games \\
\hline

\multirow{1}{*}{Interact} &
With the virtual environment (Objects, NPCs, Animals, Items, Plants, etc.) &
The Sims, Minecraft, Stray, Arida, Stardew Valley, The Stanley Parable, Among Us \\
\hline

\multirow{2}{*}{Listen and Read} &
Story/Narratives produced by the game &
The Stanley Parable, Stray, Stardew Valley, Arida \\

&
Conversations with other players &
Gartic Phone, Spyfall, Among Us, Minecraft \\
\hline

\multirow{5}{*}{Create} &
(Orally) Questions and Answers to other players &
Spyfall, Among Us, Minecraft \\

&
Phrases &
Gartic Phone, Spyfall, Minecraft \\

&
Drawings &
Gartic Phone \\

&
Objects/New environments &
The Sims, Minecraft, Stardew Valley \\

&
Characters &
The Sims, Stardew Valley, Minecraft, Among Us \\
\hline

\multirow{2}{*}{Perform Tasks} &
Created by the game &
The Sims, Among Us, The Stanley Parable, Stray, Stardew Valley, Arida \\

&
Defined by players &
Gartic Phone, Spyfall, Minecraft \\
\hline

\end{tabular}

\vspace{0.3cm}

\footnotesize
\source{The author.}
\end{threeparttable}
\end{table}

The first category, \textbf{Interact}, refers to the player’s ability to engage with the virtual environment, including objects, non-player characters (NPCs), animals, and other elements within the game. Interaction is a broad category that encompasses actions such as touching, grabbing, attacking, exploring, and conversing with NPCs. As an example, in Minecraft, players interact with the environment by building structures, mining resources, and collaborating with other players. This category highlights the importance of the virtual environment as a space for language practice, where students can engage in contextualized communication and problem-solving.

In alignment with the affordances found in the \textbf{Interact} category, the participants Manacá and Girassol suggest educational activities using the game The Sims. According to their proposals, students can use the game to expand their vocabulary by identifying familiar expressions, such as “wake up” and “go to school,” while also discovering new ones provided by the game itself. The participants highlighted how the game gives them the freedom to reproduce various actions similar to everyday life, such as walking, studying, talking, and reacting to conversations. Manacá proposed an activity in which the students create a sequence of images (by capturing the computer screen and making picture presentations) of their routines, illustrated by activities carried out in-game. This points to an integration of knowledge development through a previously defined objective.

\textbf{Listen and Read} focuses on the player’s engagement with the game’s narrative, instructions, and in-game texts. It includes listening to dialogues, reading tutorials, and following storylines to progress in the game. For instance, in The Stanley Parable, players must listen to the narrator and read on-screen instructions to navigate the game’s branching storylines. This affordance is particularly useful for developing listening and reading comprehension skills, as players are exposed to authentic language use within a meaningful context.

During the third week of the course, which focused on single-player games, participants created lesson plans that integrated reading or listening to narrative elements as a means of story progression. They also made activity lists that correlated the affordances of \textbf{Listening and Reading} with game progression. The activities highlighted that to complete the tasks and progress through the game’s narrative, students must immerse themselves in a multimodal reading experience, which simultaneously fosters language learning. Playing the games presents students with multimodal textual clues, essential for answering the questions from the proposed activities. Consequently, the activities inherently rely on students’ own ability to \textbf{listen} to \textbf{and read} the stories embedded within the games.

The \textbf{Create} category encompasses the player’s ability to create or modify elements within the game, such as drawing, writing, building, or customizing characters and environments. For example, in Gartic Phone, players produce sentences and drawings based on prompts, which can then be used to practice vocabulary and grammar. In The Sims, players create characters and environments, allowing for creative expression and language practice related to descriptions and storytelling. Furthermore, the production of questions, sentences, and drawings extends engagement beyond the game’s internal mechanics, maintaining active participant involvement throughout the learning process.

Manacá, Lavanda, and Sálvia leveraged the customization features of characters and environments in The Sims to design a lesson in which students use the game to engage with content related to the class. Afterward, they transform their “screenshots” into virtual presentations. Manacá’s proposal involved the creation of a storyboard, a sequence of images accompanied by captions that depict each student’s daily routines. The integration of the game into additional language practice provides context to the activity and offers students creative opportunities for autonomous learning, allowing them to engage with the content at home.

All participants’ creations align with previous research concerning The Sims and the exploration of multimodality in the classroom, utilizing tools for capturing screens, recording, and editing images \cite{purushotma2006sims,valadares2022videogames}.

Gérbera planned an activity of making voiceovers for the actions from the game, which affords sentence creation. The participant established connections between writing and speaking practices, allowing the game to be explored to its fullest potential. The Sims’ limitation, its absence of voiceovers, became a teaching pathway for the participant. The task of creating voiceovers draws on various multimodal resources that extend beyond the game itself. To complete the activity, students are required to write a script, plan characters, capture gameplay videos, and perform voice acting. This transforms the lesson plan into a long-term project, enhancing opportunities for practicing the additional language through diverse interactions.

The final category, \textbf{Perform Tasks}, refers to the completion of pre-programmed tasks (commonly referred to as quests) within the game, which are often presented as objectives that players must achieve to progress. These tasks can be created by the game itself or by other players. For example, in Stardew Valley, players complete quests such as planting crops or fishing, which require them to follow instructions and engage with the game’s mechanics. This affordance is useful for practicing imperative language and following sequential instructions, as well as for fostering goal-oriented learning. Manacá developed an activity centered around the quests that are present in Stardew Valley, highlighting opportunities for linguistic knowledge construction available within the game that can be revisited in a list of activities.

The activity list produced by Manacá leverages the completion of tasks, as the questions and phrases are connected to the progression of the story within the digital game. One question, “What is the first main goal of the game after these changes?”, clearly exemplifies the relationship between the game’s missions and the proposed activity. The inclusion of screenshots on the activity sheet adds depth to the participant’s planning, as they depict scenes of player interaction with the game environment, thereby reinforcing the presence of this affordance in the teaching of an additional language.


\section{Discussion and conclusion}\label{sec-formato}
The primary objective of this research was to analyze the affordances perceived by additional language teachers when using digital games for teaching and learning. This was achieved through the design and implementation of the “Play, Design \& Teach” workshop, which aimed to stimulate the production of game-based activities aligned with Ludic Language Pedagogy \cite{york2021playing}. The research sought to answer the following question: \textit{What affordances are perceived by additional language teachers in training and in service when planning and producing teaching materials that explore the use of digital games?}

The workshop revealed that participants viewed digital games as a viable tool for promoting student autonomy and engagement. The teaching materials produced during the course demonstrated a shift from traditional vocabulary-focused activities to more holistic, multimodal learning experiences. For instance, activities in Minecraft and The Sims allowed students to engage with language through creative production and interaction with virtual environments.

According to \textcite{york2021playing}, the integration of games in teaching is not confined to what is present within the game itself but can also occur in relation to its surrounding context or ecological framework. In this regard, the participating teachers proposed activities that utilized digital games not solely linked to situations within the game, but also to the culture and community surrounding them, aiming to stimulate students’ critical thinking. Such awareness can be connected to the Epistemological Curiosity defined by \textcite{freire2011autonomia}, which advocates for a change in the educational paradigm that emphasizes the necessity of “thinking rightly” in response to the challenges of teaching, considering the student as an integral part of the learning process.

Regarding the teaching materials produced for additional language instruction, participating teachers sought to integrate digital games into various activities, which, in many cases, transcended the mere “acquisition of vocabulary” emphasized in their initial narratives. Additionally, the teaching materials aimed to provide students with opportunities to construct their own knowledge in the additional language, facilitated, or guided, by digital games. The activities designed games and tasks so closely intertwined that the game and the activities became inseparable, offering opportunities for a multimodal learning experience that, in some instances, surpassed the media boundaries of digital games.

In this context, we underscore the distinction between an integrative methodology of games, as outlined by LLP, and gamification. We recognize that the concept of gamification proposes a shift in the teaching paradigm, with a primary focus on motivation. However, its reality is presented by mostly extrinsic motivation driven by points, badges, or leaderboards, which does not integrate into innovative teaching practices \cite{leffa2020gamificacao,braga2020games}. Co-opted by the educational and business market in its relentless pursuit of efficiency, \textit{gamification} has become a common presence in educational and work environments. Consequently, its mechanical and immediate approach, emphasizing mere knowledge transfer, aligns with the ‘banking education’ model \cite{freire2013oprimido} and lacks commitment to meaningful social transformation.

This contrast becomes clear upon analyzing the tasks created by the participants and the student’s role they consequently shape: while a gamified approach might position the student as a \textit{competitor} merely earning points on a quiz, the LLP-inspired projects, such as Gérbera’s voiceovers or Manacá’s storyboards, cast the student as a co-creator, providing genuine immersion to bring their stories to life.

The affordances of digital games for learning and teaching additional languages perceived by teachers were categorized into four main areas: \textbf{Interact}, \textbf{Listen and Read}, \textbf{Create}, and \textbf{Perform Tasks}. In the Interact category, we highlighted the use of the virtual environment as part of the teaching activity through interaction with various elements present within it. Activities in the \textbf{Listen and Read} category benefit from the game’s narrative and story, or from conversations among players aimed at achieving a common goal. The \textbf{Create} category encompasses a broader range of subcategories, reflecting the different ways games engage player production that can be utilized in teaching. This production can range from small insertions of written or oral phrases to more impactful creations within the virtual space, such as drawings, objects, environments, and characters. Finally, the \textbf{Perform Tasks} affordance presents situations where the player (learner) must progress according to certain rules, which may be defined by the game or by other players. Activities that draw on this category are integrated with game progression and are often associated with other affordances, such as \textbf{Interact} or \textbf{Listen and Read}.

In light of the responses shown in this research, we can assert that the use of digital games in teaching proves fruitful in its possibilities, offering teachers and learners new opportunities to connect with additional languages. To achieve this, it is essential that teachers possess an epistemological curiosity \cite{freire2011autonomia}, allowing themselves to be challenged to rethink their practices and their role in teaching while recognizing the autonomy and individuality of each student in relation to their realities. While challenges remain, the potential of digital games to transform language education is undeniable, offering new opportunities for both teachers and learners.

It is important to acknowledge the limitations of this study, which, due to its qualitative and exploratory nature, focused on a small group of six teachers within a specific professional development context. The findings presented here are not intended for generalization but rather to offer an in-depth look at the perceptions and practices that emerged in this setting. Subsequent research that follows teachers after their participation in similar workshops could analyze the long-term sustainability of these practices in the classroom. It would also be valuable to investigate the perspective of the learners themselves by analyzing their engagement and the learning outcomes resulting from game-based pedagogical activities, thereby enriching the field with a more complete view of the phenomenon.

\begingroup
\linespread{0.95}\selectfont
\printbibliography\label{sec-bib}
\endgroup
% if the text is not in Portuguese, it might be necessary to use the code below instead to print the correct ABNT abbreviations [s.n.], [s.l.]
%\begin{portuguese}
%\printbibliography[title={Bibliography}]
%\end{portuguese}


\begin{dataavailability}
\txtdataavailability{dataavailable} % options: dataavailable, dataonly, databody, datanotav, nodata
\end{dataavailability}

\begin{funding}
This study was financed by the Coordenação de Aperfeiçoamento de Pessoal de Nível Superior – Brazil (CAPES).
\end{funding}




\end{document}

